\chapter{Learning the Parameters: EM on the Chain}\label{ch:em-method}

Everything so far assumed the transition kernel was given. This chapter drops
that assumption. The data are still a Markov chain observed through the
Ornstein--Uhlenbeck channel, but the kernel that generated them is unknown and
must be recovered from the noised sequences alone: no clean data, no
supervision, and no differentiation through the inference.

The tool is expectation--maximisation, and the reason it is cheap here is
Fisher's identity. At each current parameter value, one complete
forward--backward pass of belief propagation already contains the exact
gradient of the marginal log-likelihood with respect to the parameters, so
each E-step is the BP recursion of Chapter~\ref{ch:gaussian-bp} with the
kernel held fixed and the M-step moves the kernel; the full fit runs one such
pass per iteration, and nothing is ever back-propagated through the recursion.
Section~\ref{sec:em-algorithm} derives that claim in full, in the form
actually implemented.

\section{Model: an estimated mixture innovation}\label{sec:em-model}

The clean chain is the one of \eqref{eq:l-chain-prior}, in the index
convention of this chapter,
\begin{equation}
  a_1 \sim \Nd(0,1), \qquad a_i = \alpha\, a_{i-1} + \varepsilon_i,
  \qquad i = 2,\dots,\nsites,
  \label{eq:em-chain}
\end{equation}
with innovations $\varepsilon_i$ i.i.d., mean zero and variance
$\ivar = 1-\alpha^2$: the same Gaussian initial state and the same
variance-matching Chapter~\ref{ch:laplace} used, but now drawn from a member of
a \emph{parametric family} rather than a fixed law. Forward noising is the
coordinatewise Ornstein--Uhlenbeck channel of Section~\ref{sec:ou},
$x = e^{-t} a + \sqrt{\Deltat}\,\xi$, unchanged from every earlier chapter.

The estimated family is a $C$-component Gaussian mixture on the innovation. It
is the transition kernel written $\kernel(a' \mid a)$ in
Chapters~\ref{ch:background} and~\ref{ch:laplace}, carrying its parameters
explicitly because they are now what is being estimated:
\begin{equation}
  \kernel_\theta(a' \mid a) = \sum_{c=1}^{C} \pi_c \,
    \Nd\!\big(a' - \alpha a;\ \nu_c,\ s_c^2\big),
  \qquad \theta = \{\alpha, (\pi_c,\nu_c,s_c)_{c=1}^C\},
  \label{eq:em-mixture-kernel}
\end{equation}
so both the autoregressive coefficient $\alpha$ and the shape of the innovation
density are free, estimated jointly by maximum marginal likelihood from noised
data alone. The parameter space is
\begin{equation}
  \Theta = \Big\{\, \pi_c \geq 0,\ \textstyle\sum_c \pi_c = 1,\
    s_c^2 \geq s_{\min}^2,\ \nu_c \in \R,\
    \alpha \in (-1, 1) \,\Big\},
  \label{eq:em-parameter-space}
\end{equation}
the variance floor $s_{\min}^2$ being the guard of
Section~\ref{sec:em-numerical} against a component collapsing onto a single
grid point, and $|\alpha| < 1$ being what gives the recursion a stationary law
at all. It buys existence and no more: with $(\pi_c, \nu_c, s_c)$ free the
fitted chain's stationary variance is
$\sum_c \pi_c(\nu_c^2 + s_c^2)/(1-\alpha^2)$, which equals one only under a
variance-matching constraint this family deliberately does not impose, and its
stationary mean is zero only if $\sum_c \pi_c \nu_c = 0$, which it likewise
does not impose --- the asymmetry recorded below. The conditional maximiser for
$\alpha$ derived below is a closed form that does not enforce its own
constraint, so the boundary is a real possibility rather than an excluded one;
in every fit reported in this thesis the estimate stays far inside the
interval, near the generating value, and the constraint never binds. Component
labels are exchangeable throughout, so $\Theta$ is identified only up to
permutation and every label-dependent diagnostic is reported under a fixed
ordering convention (Section~\ref{sec:em-diagnostics}). At $C=1$ this collapses to a single Gaussian, recovering the model
of Chapter~\ref{ch:gaussian}; larger $C$ can represent skew, heavy tails and
multimodality to the resolution of the grid.

One deliberate asymmetry between the generating and the fitted model should be
on the record here rather than discovered later. The \emph{generating}
innovations are centred with variance $1-\alpha^2$; the \emph{fitted} family
\eqref{eq:em-mixture-kernel} enforces neither constraint --- no
$\sum_c \pi_c\nu_c = 0$ and no matched total variance. That is a choice, not
an oversight: projecting the means back to zero after each update does not
maximise the surrogate and measurably broke monotone ascent when tried
(by $\sim 10^{-2}$ nats per edge), whereas the unconstrained ECM is exactly
monotone, and on centred data the fitted total innovation mean lands at the
$10^{-3}$ level on its own. The moment conditions are therefore treated as
\emph{diagnostics} of the fit, reported rather than imposed, and the
free-parameter count of the family is $(C-1) + C + C + 1 = 3C$. 


\section{The EM algorithm, derived}\label{sec:em-algorithm}

The parameters of the transition kernel are not observed directly. What we
observe is the noised sequence
\[
  x=(x_1,\ldots,x_{\nsites}),
\]
while the clean trajectory
\[
  a=(a_1,\ldots,a_{\nsites})
\]
that generated it is latent. In the Gaussian-mixture model there is a second
source of missing information: for every transition
$a_{i-1}\to a_i$, we do not observe which mixture component generated the
innovation. Denote by $c_i\in\{1,\ldots,C\}$ the label of the transition into $a_i$ and
write
\[
  c=(c_2,\ldots,c_{\nsites}),
\]
one label per transition, so $\nsites-1$ of them rather than $\nsites$: the
first state is drawn from the initial law and no innovation generates it.

The statistical goal is still ordinary maximum likelihood:
\begin{equation}
  \widehat\theta
  \in
  \operatorname{argmax}_\theta \mathcal L(\theta),
  \qquad
  \mathcal L(\theta)
  :=
  \log p_\theta(x).
  \label{eq:em-objective}
\end{equation}
The difficulty is that $p_\theta(x)$ is obtained only after removing all the
unobserved variables,
\begin{equation}
  p_\theta(x)
  =
  \sum_c
  \int_{\R^{\nsites}}
  p_\theta(a,c,x)\,\dd a,
  \label{eq:em-marginalisation}
\end{equation}
and hence
\begin{equation}
  \mathcal L(\theta)
  =
  \log
  \left[
    \sum_c
    \int_{\R^{\nsites}}
    p_\theta(a,c,x)\,\dd a
  \right].
  \label{eq:em-marginal-likelihood}
\end{equation}
The logarithm therefore sits \emph{outside} an integral over the whole clean
trajectory and a sum over all mixture-label configurations. Directly
optimising this expression is substantially harder than fitting the model if
$a$ and $c$ were known.

Expectation--maximisation (EM) exploits exactly this distinction. At iteration
$k$, it alternates between two questions:
\begin{enumerate}
  \item[\emph{E-step.}]
  Given the current parameters $\theta^{(k)}$ and the observation $x$, which
  hidden trajectories and mixture labels are plausible?
  \item[\emph{M-step.}]
  If those possibilities are weighted according to their current posterior
  probabilities, which new parameters fit them best?
\end{enumerate}
Thus EM does not replace the hidden variables by a single reconstruction.
It averages over their entire posterior uncertainty.

\subsection{Complete data and the surrogate}

To formalise this idea, collect all latent variables into
\begin{equation}
  z=(a,c)
\end{equation}
and define the posterior under the current parameter value,
\begin{equation}
  q^{(k)}(a,c)
  :=
  p_{\theta^{(k)}}(a,c\mid x).
  \label{eq:em-q-def}
\end{equation}
The distribution $q^{(k)}$ is computed in the E-step and is then held fixed
while the parameters are changed in the M-step.

The central EM quantity is
\begin{equation}
  Q(\theta\mid\theta^{(k)})
  :=
  \E_{q^{(k)}}\!
  \left[
    \log p_\theta(a,c,x)
  \right].
  \label{eq:em-Q-def}
\end{equation}
This is called the EM \emph{surrogate}. Its meaning is: imagine that
$(a,c)$ were observed. Then $\log p_\theta(a,c,x)$ would be the ordinary
complete-data log-likelihood. Since $(a,c)$ are not observed, EM averages that
quantity over all their possible values, weighted by how plausible those
values are under the current model.

The connection between this surrogate and the true likelihood can be derived
exactly. Bayes' rule gives
\begin{equation}
  p_\theta(a,c\mid x)
  =
  \frac{p_\theta(a,c,x)}{p_\theta(x)},
\end{equation}
and therefore
\begin{equation}
  \log p_\theta(x)
  =
  \log p_\theta(a,c,x)
  -
  \log p_\theta(a,c\mid x).
  \label{eq:em-bayes-log}
\end{equation}
Since the left-hand side does not depend on $(a,c)$, taking expectation with
respect to the fixed distribution $q^{(k)}$ leaves it unchanged:
\begin{align}
  \mathcal L(\theta)
  &=
  \E_{q^{(k)}}\!
  \left[
    \log p_\theta(a,c,x)
  \right]
  -
  \E_{q^{(k)}}\!
  \left[
    \log p_\theta(a,c\mid x)
  \right]
  \nonumber\\
  &=
  Q(\theta\mid\theta^{(k)})
  -
  \E_{q^{(k)}}\!
  \left[
    \log p_\theta(a,c\mid x)
  \right].
  \label{eq:em-first-decomposition}
\end{align}

Now add and subtract
$\E_{q^{(k)}}[\log q^{(k)}(a,c)]$. Using
\begin{equation}
  \mathcal H(q)
  :=
  -\E_q[\log q]
\end{equation}
and
\begin{equation}
  \operatorname{KL}(q\,\|\,p)
  :=
  \E_q\!\left[
    \log\frac{q}{p}
  \right],
\end{equation}
we obtain
\begin{equation}
  \boxed{
  \mathcal L(\theta)
  =
  Q(\theta\mid\theta^{(k)})
  +
  \mathcal H(q^{(k)})
  +
  \operatorname{KL}\!\left(
    q^{(k)}
    \,\middle\|\,
    p_\theta(a,c\mid x)
  \right).
  }
  \label{eq:em-decomposition}
\end{equation}
Equivalently,
\begin{equation}
  \underbrace{
    Q(\theta\mid\theta^{(k)})
    +
    \mathcal H(q^{(k)})
  }_{\displaystyle \mathcal F(\theta;q^{(k)})}
  =
  \mathcal L(\theta)
  -
  \operatorname{KL}\!\left(
    q^{(k)}
    \,\middle\|\,
    p_\theta(a,c\mid x)
  \right)
  \leq
  \mathcal L(\theta).
  \label{eq:em-lower-bound}
\end{equation}
Hence
$\mathcal F(\theta;q^{(k)})$ is a lower bound on the true observed-data
log-likelihood. During the M-step $q^{(k)}$ is fixed, so its entropy is a
constant with respect to $\theta$. Maximising the lower bound
$\mathcal F$ is therefore exactly equivalent to maximising $Q$.

This explains the word \emph{surrogate}: rather than attacking
$\mathcal L(\theta)$ directly, EM optimises the easier function
$Q(\theta\mid\theta^{(k)})$ while keeping the current posterior fixed.

There is one further crucial property. At the current parameter
$\theta=\theta^{(k)}$,
\begin{equation}
  q^{(k)}(a,c)
  =
  p_{\theta^{(k)}}(a,c\mid x),
\end{equation}
so
\begin{equation}
  \operatorname{KL}\!\left(
    q^{(k)}
    \,\middle\|\,
    p_{\theta^{(k)}}(a,c\mid x)
  \right)
  =0.
\end{equation}
Consequently the lower bound touches the true likelihood at the current
iterate,
\begin{equation}
  \mathcal F(\theta^{(k)};q^{(k)})
  =
  \mathcal L(\theta^{(k)}).
  \label{eq:em-touching}
\end{equation}
The EM iteration can therefore be pictured as
\[
  \theta^{(k)}
  \;\longrightarrow\;
  q^{(k)}(a,c)
  \;\longrightarrow\;
  Q(\theta\mid\theta^{(k)})
  \;\longrightarrow\;
  \theta^{(k+1)}.
\]
The first arrow infers what the hidden clean process probably was; the second
turns that posterior into a weighted complete-data likelihood; the final arrow
fits a new transition kernel to those posterior-weighted hidden trajectories.

In a generic latent-variable model, computing $q^{(k)}$ may itself require an
approximation. Here the chain structure makes this step particularly simple.

\begin{proposition}[The E-step is exact]\label{prop:em-estep}
Since the posterior factor graph is a tree, one forward and one backward
sum--product sweep return the exact single-site and pairwise posterior
marginals, at the level of functional messages. Consequently the E-step
requires neither a variational approximation nor Monte Carlo sampling: all
conditional expectations needed by EM are exact under the model, before the
numerical grid discretisation is introduced.
\end{proposition}

We next write the complete-data likelihood explicitly. Introduce the mixture
label $c_i$ on the transition $a_{i-1}\to a_i$. Conditional on $c_i=c$, the
mixture transition becomes a single Gaussian,
\begin{equation}
  \widetilde\kernel_\theta(a',c\mid a)
  =
  \pi_c\,
  \Nd\!\big(
    a'-\alpha a;
    \nu_c,s_c^2
  \big),
  \qquad
  \kernel_\theta(a'\mid a)
  =
  \sum_{c=1}^{C}
  \widetilde\kernel_\theta(a',c\mid a).
  \label{eq:em-complete-kernel}
\end{equation}

The purpose of introducing $c$ is now visible. Without the label, one
transition contributes
\begin{equation}
  \log\kernel_\theta(a_i\mid a_{i-1})
  =
  \log
  \left[
    \sum_{c=1}^{C}
    \pi_c
    \Nd\!\big(
      a_i-\alpha a_{i-1};
      \nu_c,s_c^2
    \big)
  \right],
  \label{eq:em-log-mixture}
\end{equation}
so the logarithm couples all mixture components through a log--sum. Once
$c_i$ is exposed, the same transition is assigned to a single Gaussian
component, giving a weighted Gaussian likelihood with closed-form conditional
maximisers.

For one sequence, the complete joint distribution factorises as
\begin{equation}
  p_\theta(a,c,x)
  =
  p(a_1)
  \prod_{i=2}^{\nsites}
  \widetilde\kernel_\theta(a_i,c_i\mid a_{i-1})
  \prod_{i=1}^{\nsites}
  \phi_i(a_i),
  \label{eq:em-complete-factorisation}
\end{equation}
where
\[
  \phi_i(a_i)=p(x_i\mid a_i)
\]
is the fixed OU observation factor. Taking logarithms,
\begin{equation}
  \log p_\theta(a,c,x)
  =
  \underbrace{\log p(a_1)}_{\text{independent of }\theta}
  +
  \sum_{i=2}^{\nsites}
  \log
  \widetilde\kernel_\theta(a_i,c_i\mid a_{i-1})
  +
  \underbrace{
    \sum_{i=1}^{\nsites}\log\phi_i(a_i)
  }_{\text{independent of }\theta}.
  \label{eq:em-complete-ll}
\end{equation}
Thus, under the parameterisation used here, all terms relevant to the
M-step are local transition terms.

For one labelled transition,
\begin{align}
  \log
  \widetilde\kernel_\theta(a_i,c_i\mid a_{i-1})
  &=
  \log\pi_{c_i}
  +
  \log
  \Nd\!\big(
    a_i-\alpha a_{i-1};
    \nu_{c_i},s_{c_i}^2
  \big)
  \nonumber\\
  &=
  \log\pi_{c_i}
  -\frac12\log(2\pi)
  -\frac12\log s_{c_i}^2
  -
  \frac{
    (a_i-\alpha a_{i-1}-\nu_{c_i})^2
  }{2s_{c_i}^2}.
  \label{eq:em-single-transition-log}
\end{align}
Using indicators,
\begin{align}
  \log
  \widetilde\kernel_\theta(a_i,c_i\mid a_{i-1})
  =
  \sum_{c=1}^{C}
  \mathbf 1\{c_i=c\}
  \left[
    \log\pi_c
    -\frac12\log(2\pi)
    -\frac12\log s_c^2
    -
    \frac{
      (a_i-\alpha a_{i-1}-\nu_c)^2
    }{2s_c^2}
  \right].
  \label{eq:em-indicator-transition}
\end{align}

Substituting this expression into \eqref{eq:em-Q-def} gives
\begin{align}
  Q(\theta\mid\theta^{(k)})
  &=
  \E_{a,c\mid x,\theta^{(k)}}
  \!\left[
    \log p_\theta(a,c,x)
  \right]
  \nonumber\\
  &\overset{c}{=}
  \sum_{i=2}^{\nsites}
  \sum_{c=1}^{C}
  \E_{\theta^{(k)}}\!
  \left[
    \mathbf 1\{c_i=c\}
    \left(
      \log\pi_c
      -\frac12\log s_c^2
      -
      \frac{
        (a_i-\alpha a_{i-1}-\nu_c)^2
      }{2s_c^2}
    \right)
  \right],
  \label{eq:em-Q}
\end{align}
where $\overset{c}{=}$ means equality up to an additive constant independent
of $\theta$. In particular, the initial-state term, the observation-channel
term and $-\frac12\log(2\pi)$ do not affect the M-step.

Equation~\eqref{eq:em-Q} is the central simplification: the difficult
logarithm of the marginalisation in
\eqref{eq:em-marginal-likelihood} has been replaced by the expectation of a
sum of local Gaussian transition terms. The remaining question is how those
expectations are computed.

\subsection{Fisher's identity: why BP is never differentiated}

The same construction also explains why parameter learning does not require
differentiating through the BP recursion. Starting from
\begin{equation}
  \mathcal L(\theta)
  =
  \log p_\theta(x),
\end{equation}
differentiate directly. The interchange of $\nabla_\theta$ with the sum and
the integral below is the one regularity condition the identity needs: it holds
because $p_\theta(a,c,x)$ is differentiable in $\theta$ with
$\|\nabla_\theta p_\theta(a,c,x)\|$ dominated, locally uniformly in
$\theta$, by an integrable function of $(a,c)$ --- which the Gaussian
likelihood factors and the mixture kernel of \eqref{eq:em-mixture-kernel}
satisfy on a compact subset of $\Theta$ containing every fit reported here ---
every component non-degenerate and the fitted $\alpha$ staying well inside
$(-1,1)$. Then
\begin{align}
  \nabla_\theta\mathcal L(\theta)
  &=
  \frac{1}{p_\theta(x)}
  \nabla_\theta p_\theta(x)
  \nonumber\\
  &=
  \frac{1}{p_\theta(x)}
  \sum_c
  \int
  \nabla_\theta p_\theta(a,c,x)\,\dd a.
  \label{eq:em-fisher-step1}
\end{align}
Using
\begin{equation}
  \nabla_\theta p_\theta(a,c,x)
  =
  p_\theta(a,c,x)
  \nabla_\theta\log p_\theta(a,c,x),
\end{equation}
we obtain
\begin{align}
  \nabla_\theta\mathcal L(\theta)
  &=
  \sum_c
  \int
  \frac{p_\theta(a,c,x)}{p_\theta(x)}
  \nabla_\theta
  \log p_\theta(a,c,x)\,\dd a
  \nonumber\\
  &=
  \sum_c
  \int
  p_\theta(a,c\mid x)
  \nabla_\theta
  \log p_\theta(a,c,x)\,\dd a
  \nonumber\\
  &=
  \E_{a,c\mid x,\theta}
  \left[
    \nabla_\theta
    \log p_\theta(a,c,x)
  \right].
  \label{eq:em-fisher-general}
\end{align}
This is Fisher's identity.

Since only the transition factors in
\eqref{eq:em-complete-ll} depend on $\theta$,
\begin{equation}
  \nabla_\theta\mathcal L(\theta)
  =
  \E_{a,c\mid x,\theta}
  \left[
    \sum_{i=2}^{\nsites}
    \nabla_\theta
    \log
    \widetilde\kernel_\theta
    (a_i,c_i\mid a_{i-1})
  \right].
  \label{eq:em-fisher}
\end{equation}
Evaluating at the current iterate gives
\begin{equation}
  \nabla_\theta\mathcal L(\theta)
  \Big|_{\theta=\theta^{(k)}}
  =
  \nabla_\theta
  Q(\theta\mid\theta^{(k)})
  \Big|_{\theta=\theta^{(k)}}.
  \label{eq:em-fisher-Q}
\end{equation}

The important point is what is \emph{not} differentiated. BP is used only to
compute
\[
  q^{(k)}(a,c)
  =
  p_{\theta^{(k)}}(a,c\mid x).
\]
Once that posterior has been computed, it is treated as fixed. The derivative
acts on the candidate kernel parameters inside
$\log\widetilde\kernel_\theta$, not on the BP messages that produced the
posterior. Thus BP supplies the posterior expectations and Fisher's identity
converts them into the likelihood gradient without back-propagating through
the inference algorithm.

\subsection{The E-step: which expectations BP supplies}

The complete-data objective in \eqref{eq:em-Q} is a sum over transitions.
Consequently we do not need the full joint posterior over the entire clean
trajectory explicitly. It is sufficient to know, for each edge
$i-1\to i$, the pairwise posterior
\begin{equation}
  b_i(a_{i-1},a_i)
  :=
  p_{\theta^{(k)}}(a_{i-1},a_i\mid x).
  \label{eq:em-pairwise-def}
\end{equation}
On a chain this marginal is obtained exactly from the forward and backward
messages:
\begin{equation}
  b_i(a_{i-1},a_i)
  \propto
  \lmsg_{i-1}(a_{i-1})
  \phi_{i-1}(a_{i-1})
  \kernel_{\theta^{(k)}}(a_i\mid a_{i-1})
  \phi_i(a_i)
  \rmsg_i(a_i),
  \label{eq:em-pairwise}
\end{equation}
with normalisation
\begin{equation}
  \int\!\!\int
  b_i(a_{i-1},a_i)
  \,\dd a_{i-1}\dd a_i
  =1.
\end{equation}

The mixture label requires no additional BP recursion. Once a pair
$(a_{i-1},a_i)$ is given, Bayes' rule on the finite mixture gives the
responsibility
\begin{align}
  r_{i,c}^{(k)}(a_{i-1},a_i)
  &:=
  \Prob_{\theta^{(k)}}\!
  \left(
    c_i=c
    \mid
    a_{i-1},a_i
  \right)
  \nonumber\\
  &=
  \frac{
    \pi_c^{(k)}
    \Nd\!\left(
      a_i-\alpha^{(k)}a_{i-1};
      \nu_c^{(k)},(s_c^2)^{(k)}
    \right)
  }{
    \displaystyle
    \sum_{d=1}^{C}
    \pi_d^{(k)}
    \Nd\!\left(
      a_i-\alpha^{(k)}a_{i-1};
      \nu_d^{(k)},(s_d^2)^{(k)}
    \right)
  }
  \nonumber\\
  &=
  \frac{
    \pi_c^{(k)}
    \Nd\!\left(
      a_i-\alpha^{(k)}a_{i-1};
      \nu_c^{(k)},(s_c^2)^{(k)}
    \right)
  }{
    \kernel_{\theta^{(k)}}(a_i\mid a_{i-1})
  }.
  \label{eq:em-responsibility}
\end{align}
For every pair,
\begin{equation}
  \sum_{c=1}^{C}
  r_{i,c}^{(k)}(a_{i-1},a_i)
  =1.
\end{equation}

Define expectation with respect to the pairwise posterior by
\begin{equation}
  \langle f\rangle_i
  :=
  \int\!\!\int
  f(a_{i-1},a_i)
  b_i(a_{i-1},a_i)
  \,\dd a_{i-1}\dd a_i.
  \label{eq:em-bracket}
\end{equation}
Because
\begin{equation}
  \E\!\left[
    \mathbf 1\{c_i=c\}
    \mid a_{i-1},a_i,x,\theta^{(k)}
  \right]
  =
  r_{i,c}^{(k)}(a_{i-1},a_i),
  \label{eq:em-indicator-responsibility}
\end{equation}
every expectation in \eqref{eq:em-Q} can be reduced to an integral against
$b_i$ weighted by $r_{i,c}^{(k)}$.

For each component define the six sufficient statistics
\begin{equation}
  \begin{aligned}
    R_c
    &=
    \sum_i
    \left\langle r_{i,c}^{(k)}\right\rangle_i,
    &
    U_c
    &=
    \sum_i
    \left\langle
      r_{i,c}^{(k)}a_i
    \right\rangle_i,
    &
    V_c
    &=
    \sum_i
    \left\langle
      r_{i,c}^{(k)}a_{i-1}
    \right\rangle_i,
    \\[1mm]
    P_c
    &=
    \sum_i
    \left\langle
      r_{i,c}^{(k)}a_{i-1}a_i
    \right\rangle_i,
    &
    W_c
    &=
    \sum_i
    \left\langle
      r_{i,c}^{(k)}a_{i-1}^2
    \right\rangle_i,
    &
    G_c
    &=
    \sum_i
    \left\langle
      r_{i,c}^{(k)}a_i^2
    \right\rangle_i.
  \end{aligned}
  \label{eq:em-suffstats}
\end{equation}
For several training sequences the same quantities are simply summed over
all edges of all sequences.

 $R_c$ is the posterior effective
number of transitions assigned to component $c$; $U_c$ and $V_c$ are
responsibility-weighted first moments of the child and parent states;
$G_c$ and $W_c$ are the corresponding second moments; and $P_c$ is the
parent--child cross moment.

A single forward--backward pass therefore reduces the entire E-step to these
posterior-weighted sufficient statistics.

\subsection{The M-step: deriving the conditional maximisers}

We now substitute \eqref{eq:em-suffstats} into the surrogate and derive the
parameter updates explicitly.

First expand the squared residual:
\begin{align}
  (a_i-\alpha a_{i-1}-\nu_c)^2
  &=
  a_i^2
  -2\alpha a_{i-1}a_i
  +\alpha^2 a_{i-1}^2
  -2\nu_c a_i
  +2\alpha\nu_c a_{i-1}
  +\nu_c^2.
  \label{eq:em-residual-expand}
\end{align}
Taking the posterior-weighted sums gives
\begin{equation}
  S_c(\alpha,\nu_c)
  :=
  G_c
  -2\alpha P_c
  +\alpha^2W_c
  -2\nu_c U_c
  +2\alpha\nu_c V_c
  +\nu_c^2R_c.
  \label{eq:em-residual-stat}
\end{equation}
Hence, up to constants independent of $\theta$,
\begin{equation}
  \boxed{
  Q(\theta\mid\theta^{(k)})
  \overset{c}{=}
  \sum_{c=1}^{C}
  \left[
    R_c\log\pi_c
    -
    \frac{R_c}{2}\log s_c^2
    -
    \frac{
      S_c(\alpha,\nu_c)
    }{2s_c^2}
  \right].
  }
  \label{eq:em-Q-suffstats}
\end{equation}
All quantities $R_c,U_c,V_c,P_c,W_c,G_c$ are fixed during the M-step: they
were computed using $\theta^{(k)}$ in the E-step. Only the candidate
parameters $\pi_c,\nu_c,s_c^2,\alpha$ are varied.

\paragraph{Mixture weights.}

The part of $Q$ depending on the weights is
\begin{equation}
  Q_\pi
  =
  \sum_{c=1}^{C}
  R_c\log\pi_c,
\end{equation}
subject to
\begin{equation}
  \sum_{c=1}^{C}\pi_c=1.
\end{equation}
Introduce a Lagrange multiplier $\lambda$:
\begin{equation}
  \mathcal J
  =
  \sum_c R_c\log\pi_c
  +
  \lambda
  \left(
    \sum_c\pi_c-1
  \right).
\end{equation}
Differentiating,
\begin{equation}
  \frac{\partial\mathcal J}{\partial\pi_c}
  =
  \frac{R_c}{\pi_c}
  +
  \lambda
  =
  0,
\end{equation}
so
\begin{equation}
  \pi_c
  =
  -\frac{R_c}{\lambda}.
\end{equation}
Summing over $c$ and imposing $\sum_c\pi_c=1$ gives
\begin{equation}
  \lambda
  =
  -\sum_{d=1}^{C}R_d.
\end{equation}
Therefore
\begin{equation}
  \boxed{
  \pi_c^{+}
  =
  \frac{R_c}{
    \displaystyle\sum_{d=1}^{C}R_d
  }.
  }
  \label{eq:em-mstep-pi}
\end{equation}
Thus the new mixture weight is simply the fraction of the total posterior
transition mass assigned to component $c$.

\paragraph{Component means.}

Holding $\alpha$ and $s_c^2$ fixed, only
$S_c(\alpha,\nu_c)$ depends on $\nu_c$. From
\eqref{eq:em-residual-stat},
\begin{equation}
  \frac{\partial S_c}{\partial\nu_c}
  =
  -2U_c
  +2\alpha V_c
  +2\nu_cR_c.
\end{equation}
Since
\begin{equation}
  \frac{\partial Q}{\partial\nu_c}
  =
  -\frac{1}{2s_c^2}
  \frac{\partial S_c}{\partial\nu_c},
\end{equation}
the stationarity condition is
\begin{equation}
  -2U_c
  +2\alpha V_c
  +2\nu_cR_c
  =
  0.
\end{equation}
Hence
\begin{equation}
  \boxed{
  \nu_c^{+}
  =
  \frac{
    U_c-\alpha V_c
  }{R_c}.
  }
  \label{eq:em-mstep-nu}
\end{equation}

\paragraph{Component variances.}

Let
\[
  \tau_c=s_c^2.
\]
For fixed $\alpha$ and $\nu_c$, the relevant part of the surrogate is
\begin{equation}
  Q_c(\tau_c)
  =
  -\frac{R_c}{2}\log\tau_c
  -
  \frac{S_c(\alpha,\nu_c)}{2\tau_c}.
\end{equation}
Differentiating,
\begin{equation}
  \frac{\partial Q_c}{\partial\tau_c}
  =
  -\frac{R_c}{2\tau_c}
  +
  \frac{
    S_c(\alpha,\nu_c)
  }{2\tau_c^2}.
\end{equation}
Setting this equal to zero gives
\begin{equation}
  R_c\tau_c
  =
  S_c(\alpha,\nu_c),
\end{equation}
and therefore
\begin{equation}
  (s_c^2)^{+}
  =
  \frac{
    S_c(\alpha,\nu_c^{+})
  }{R_c}.
  \label{eq:em-mstep-var-general}
\end{equation}
Substituting
$\nu_c^{+}=(U_c-\alpha V_c)/R_c$ simplifies the weighted residual to
\begin{equation}
  S_c(\alpha,\nu_c^{+})
  =
  G_c
  -2\alpha P_c
  +\alpha^2W_c
  -
  R_c(\nu_c^{+})^2,
\end{equation}
so
\begin{equation}
  \boxed{
  (s_c^2)^{+}
  =
  \frac{
    G_c
    -2\alpha P_c
    +\alpha^2W_c
    -R_c(\nu_c^{+})^2
  }{R_c}.
  }
  \label{eq:em-mstep-var}
\end{equation}

Combining the three mixture updates,
\begin{equation}
  \boxed{
  \pi_c^{+}
  =
  \frac{R_c}{\sum_d R_d},
  \qquad
  \nu_c^{+}
  =
  \frac{U_c-\alpha V_c}{R_c},
  \qquad
  (s_c^2)^{+}
  =
  \frac{
    G_c-2\alpha P_c+\alpha^2W_c
    -R_c(\nu_c^{+})^2
  }{R_c}.
  }
  \label{eq:em-mstep-mixture}
\end{equation}

\begin{remark}[Floors, empty components, and label order]
\label{rmk:em-degeneracy}
The updates \eqref{eq:em-mstep-mixture} are undefined at $R_c = 0$ and the
mixture likelihood is unbounded as any $s_c^2 \to 0$ with a component pinned
on one residual, so the implementation carries two explicit guards: the
responsibility mass is floored at $10^{-300}$ before any division (an empty
component keeps its parameters rather than producing indeterminate ones), and
the variance update is floored at $s_c^2 \geq 10^{-6}$, a component standard
deviation of $10^{-3}$ --- far below the grid spacing, so the floor marks
degeneracy rather than shaping any resolved fit. A separate reporting gate
(\provenanceappendix) excludes fits whose narrowest component falls under two
grid cells, where the quadrature likelihood can reward lattice artefacts.
Component labels are exchangeable, so all label-dependent diagnostics order
components deterministically before comparison; the effective component count
$\exp[\mathrm{entropy}(\pi)]$ is logged per fit so that collapse onto fewer
live components is visible rather than silent.
\end{remark}

\paragraph{Autoregressive coefficient.}

Finally hold the mixture parameters fixed and optimise with respect to
$\alpha$. From \eqref{eq:em-residual-stat},
\begin{equation}
  \frac{\partial S_c}{\partial\alpha}
  =
  -2P_c
  +2\alpha W_c
  +2\nu_cV_c.
\end{equation}
Therefore
\begin{align}
  \frac{\partial Q}{\partial\alpha}
  &=
  -\sum_{c=1}^{C}
  \frac{1}{2s_c^2}
  \frac{\partial S_c}{\partial\alpha}
  \nonumber\\
  &=
  \sum_{c=1}^{C}
  \frac{
    P_c-\alpha W_c-\nu_cV_c
  }{s_c^2}.
  \label{eq:em-alpha-derivative}
\end{align}
Setting the derivative to zero gives
\begin{equation}
  \sum_c
  s_c^{-2}
  \big(
    P_c-\nu_cV_c
  \big)
  =
  \alpha
  \sum_c
  s_c^{-2}W_c,
\end{equation}
and hence
\begin{equation}
  \boxed{
  \alpha^{+}
  =
  \frac{
    \displaystyle
    \sum_c
    s_c^{-2}
    \big(
      P_c-\nu_cV_c
    \big)
  }{
    \displaystyle
    \sum_c
    s_c^{-2}W_c
  }.
  }
  \label{eq:em-mstep-alpha}
\end{equation}
The factors $s_c^{-2}$ are precisions: components with smaller innovation
variance provide more information about the regression coefficient and
therefore receive greater weight.

At $C=1$ and $\nu_1=0$ this reduces to
\begin{equation}
  \alpha^{+}
  =
  \frac{P_1}{W_1}
  =
  \frac{
    \sum_i
    \E[a_{i-1}a_i\mid x]
  }{
    \sum_i
    \E[a_{i-1}^2\mid x]
  },
\end{equation}
which is the posterior analogue of the ordinary least-squares slope of $a_i$
on $a_{i-1}$.

The complete EM iteration can therefore be summarised as
\begin{equation}
  \boxed{
  \begin{array}{ll}
  \text{E-step:}
  &
  \theta^{(k)}
  \xrightarrow{\mathrm{BP}}
  \left\{
    R_c,U_c,V_c,P_c,W_c,G_c
  \right\}_{c=1}^{C},
  \\[2mm]
  \text{M-step:}
  &
  \left\{
    R_c,U_c,V_c,P_c,W_c,G_c
  \right\}_{c=1}^{C}
  \xrightarrow{\text{closed-form updates}}
  \theta^{(k+1)}.
  \end{array}
  }
  \label{eq:em-summary}
\end{equation}

\subsection{Generalised EM and monotonicity}

The updates above are conditional maximisers. The mixture block
$(\pi,\nu,s^2)$ is maximised while $\alpha$ is fixed, and then $\alpha$ is
maximised while the mixture block is fixed. Alternating these operations
increases the same surrogate $Q$, but need not find its joint maximiser in one
step. The implementation therefore belongs to the generalised EM family, more
specifically to expectation--conditional maximisation (ECM)
\citep{meng1993maximum}.

The monotonicity guarantee does not require an exact joint maximisation of
$Q$. To see this, use \eqref{eq:em-decomposition} with
$q^{(k)}=p_{\theta^{(k)}}(a,c\mid x)$. At the current parameter the KL term
vanishes, so
\begin{align}
  \mathcal L(\theta)
  -
  \mathcal L(\theta^{(k)})
  &=
  Q(\theta\mid\theta^{(k)})
  -
  Q(\theta^{(k)}\mid\theta^{(k)})
  \nonumber\\
  &\quad+
  \operatorname{KL}\!\left(
    p_{\theta^{(k)}}(a,c\mid x)
    \,\middle\|\,
    p_\theta(a,c\mid x)
  \right).
  \label{eq:em-monotone}
\end{align}
The second term is non-negative. Consequently,
\begin{equation}
  Q(\theta^{(k+1)}\mid\theta^{(k)})
  \geq
  Q(\theta^{(k)}\mid\theta^{(k)})
\end{equation}
is sufficient to imply
\begin{equation}
  \boxed{
  \mathcal L(\theta^{(k+1)})
  \geq
  \mathcal L(\theta^{(k)}).
  }
  \label{eq:em-likelihood-monotone}
\end{equation}
Thus each conditional maximisation need only increase the surrogate; joint
maximisation is unnecessary for likelihood monotonicity.

In the implementation, the two parameter blocks are alternated until either
the gain in $Q$ per edge falls below the prescribed tolerance or a maximum of
sixteen inner sweeps is reached. The returned kernel records which stopping
condition occurred. One distinction keeps the guarantee honest. The derivation
above concerns the continuum model; the implementation evaluates densities on
the grid of Section~\ref{sec:em-numerical} without renormalising transition
columns, so the objective it computes --- and the one whose monotone increase
Table~\ref{tab:em-checks} certifies --- is the quadrature-defined surrogate of
that discretised factor model. It approximates the continuum likelihood to the
quadrature error measured in Table~\ref{tab:em-grid}, but the exact
monotonicity statement applies to the grid objective, not automatically to the
continuum one.

\begin{remark}[The inner budget of sixteen]\label{rmk:em-inner}
The cap of sixteen inner sweeps is chosen from the numerical behaviour rather
than from a theoretical convention. At $C=8$, $\nseq=256$ and a fixed outer
budget of $40$ iterations, increasing the inner budget from four to sixteen
sweeps moves the fitted innovation excess kurtosis from $1.68$ to $2.14$, a
$28\%$ change relative to the former value and still below the true Laplace
value $3.0$, whereas $\widehat\alpha$ changes only in the fourth decimal
place. The innovation shape therefore converges substantially more slowly than
the correlation parameter. By roughly $60$ outer iterations the outer and
inner loops largely compensate for one another and the inner budget has little
remaining effect. The choice of sixteen is consequently a practical safeguard
against interpreting incomplete inner optimisation as a failure to recover the
innovation shape.
\end{remark}


\section{The numerical representation: grid BP, again}
\label{sec:em-numerical}

As in Section~\ref{sec:l-grid}, messages are represented on $\ngrid$
equispaced points of $[-A,A]$ and updated by trapezoidal quadrature. The
pairwise belief \eqref{eq:em-pairwise} is therefore an
$\ngrid\times\ngrid$ array, and the expectations in
\eqref{eq:em-suffstats} become weighted sums. One EM iteration costs
$O(\nsites\ngrid^2)$ per sequence, with an additional factor $C$ only in the
mixture responsibilities.

There are two distinct numerical errors. \emph{Truncation} comes from replacing
$\R$ by $[-A,A]$ and is monitored through forward-message boundary mass;
because this depends on the sampled trajectory, it is reported over
$\gdchains$ chains. \emph{Quadrature} comes from the spacing
$h=2A/(\ngrid-1)$: in the interior the trapezoidal error scales as $O(h^2)$,
with the measured residual decreasing by a factor $4.00$ per halving. Near the
boundary it does not converge, because the remaining error is truncation rather
than resolution. Table~\ref{tab:em-grid} and Figure~\ref{fig:em-grid} separate
the two effects.

\begin{center}
\begin{minipage}{\linewidth}\centering
\captionof{table}[The two discretisation budgets]{Numerical error diagnostics.
\emph{Upper:} boundary mass versus domain half-width at $\ngrid=401$, reported
over $\gdchains$ chains. \emph{Lower:} column-mass residual versus grid spacing
at $A=8$. Interior error follows the expected $O(h^2)$ rate, whereas the edge
residual is truncation-dominated. The working choice is $\ngrid=401$, $A=8$.}
\label{tab:em-grid}
\small
\begin{tabular}{@{}lcccc@{}}
\toprule
\multicolumn{5}{@{}l}{\emph{truncation: boundary mass against half-width} ($\ngrid=401$)}\\
half-width $A$      & $4$ & $6$ & $8$ & $10$ \\
\midrule
max over chains     & $2.3\times10^{-3}$ & $1.6\times10^{-4}$ & $1.2\times10^{-6}$ & $3.2\times10^{-8}$ \\
$90$th percentile   & $2.7\times10^{-4}$ & $2.5\times10^{-6}$ & $1.5\times10^{-8}$ & $7.8\times10^{-11}$ \\
median              & $3.8\times10^{-5}$ & $3.7\times10^{-7}$ & $3.0\times10^{-9}$ & $1.8\times10^{-11}$ \\
max / median        & $61\times$ & $434\times$ & $392\times$ & $1796\times$ \\
\midrule
\multicolumn{5}{@{}l}{\emph{quadrature: column-mass residual against spacing} ($A=8$)}\\
$\ngrid$ \ / \ $h$  & \multicolumn{1}{c}{$201$ / $0.080$} & \multicolumn{1}{c}{$401$ / $0.040$} & \multicolumn{1}{c}{$801$ / $0.020$} & rate \\
\midrule
interior column     & $3.84\times10^{-3}$ & $9.61\times10^{-4}$ & $2.40\times10^{-4}$ & $4.00\times$ per halving \\
edge column         & $1.62\times10^{-2}$ & $1.90\times10^{-2}$ & $1.97\times10^{-2}$ & none \\
\bottomrule
\end{tabular}
\end{minipage}
\end{center}

\begin{figure}[t]\centering
\includegraphics[width=\linewidth]{fig_em_grid.pdf}
\caption[The two discretisation budgets, separated]{Numerical behaviour of the
grid representation. \emph{(a)} Boundary mass decreases rapidly with the
domain half-width $A$, but varies substantially across sampled chains.
\emph{(b)} Interior column-mass error follows the $O(h^2)$ trapezoidal rate,
whereas the edge error remains truncation-dominated. \emph{(c)} Relative score
error versus diffusion time for three resolutions; failure occurs when the
resolution condition \eqref{eq:l-resolution} is violated.}
\label{fig:em-grid}
\end{figure}

\paragraph{Which finite-grid model is fitted.}
The transition array is the continuum density
$\kernel_\theta(a'\mid a)$ evaluated on the grid and integrated with
trapezoidal weights, without normalising each parent-conditioned column.
BP is therefore exact, up to floating point, for the resulting discretised
factor model, but this is not exactly a normalised finite-state Markov kernel.
The discrepancy is measured by the column-mass residual in
Table~\ref{tab:em-grid}.

Column-wise renormalisation would remove this discrepancy, but its normaliser
$Z_\theta$ would depend on the parameters. Terms
$\nabla_\theta\log Z_\theta$ would then enter Fisher's identity and the
M-step, destroying the closed-form updates
\eqref{eq:em-mstep-mixture}--\eqref{eq:em-mstep-alpha}. The present convention
therefore fits the continuum kernel through the quadrature rule, with
truncation controlled separately.

\begin{remark}[The boundary diagnostic is empirical, not a bound]
\label{rmk:em-boundary}
Boundary mass depends on the realised trajectory, so a single chain cannot
certify truncation error. At the working configuration, the batch maximum over
$\gdchains$ chains exceeds the single-chain reading by a median factor
$\gdboundmed$ and by as much as $\gdboundhi$. This motivates reporting the
distribution across chains rather than one value. The quadrature comparison,
by contrast, depends only on $(A,\ngrid)$ and reproduces to $0.000$ relative
difference under reseeding.
\end{remark}

Finally, as in Section~\ref{sec:l-grid}, posterior-mean errors are amplified in
the score by
\[
  \frac{e^{-t}}{\Deltat}\sim\frac{1}{2t},
\]
so numerical or estimation deficiencies are most visible at small diffusion
times.

\section{Diagnostics}\label{sec:em-diagnostics}

The remaining checks answer two questions: whether the implementation is
correct, and which parameters can actually be recovered from noised data.

\subsection{Implementation and estimator checks}

The BP recursion agrees with independent references to machine precision, and
the surrogate never decreases in the logged runs.

\begin{center}
\begin{minipage}{\linewidth}\centering
\captionof{table}[Correctness checks on the EM implementation]{Correctness and
monotonicity checks at $\ngrid=401$, $A=8$.}
\label{tab:em-checks}
\small
\begin{tabular}{@{}llc@{}}
\toprule
check & reference & discrepancy \\
\midrule
posterior mean & closed-form Gaussian chain (Ch.~\ref{ch:gaussian})
  & $9.2\times10^{-15}$ \\
posterior mean & brute-force enumeration, no shared code
  & $1.6\times10^{-14}$ \\
log-evidence   & brute-force enumeration, no shared code
  & $1.8\times10^{-15}$ \\
\midrule
\multicolumn{2}{@{}l}{minimum increase of $Q$ over twelve steps, Gaussian kernel}
  & $+2.61$ \\
\multicolumn{2}{@{}l}{\qquad\qquad the same, Laplace kernel}
  & $+4.34$ \\
\multicolumn{2}{@{}l}{\qquad\qquad the same, $C=8$ mixture kernel}
  & $+4.18$ \\
\multicolumn{2}{@{}l}{decreases recorded in any logged run}
  & $0$ \\
\bottomrule
\end{tabular}
\end{minipage}
\end{center}

Figure~\ref{fig:em-diagnostics} adds two estimator-level checks. First,
$\widehat\alpha$ converges to the same value from initialisations spanning
$[-0.4,0.6]$, showing that it is genuinely estimated rather than supplied.
Second, its RMSE and that of the innovation variance decay approximately at the
$\nseq^{-1/2}$ sampling rate.

\begin{figure}[t]\centering
\includegraphics[width=\linewidth]{fig_em_diagnostics.pdf}
\caption[Optimisation diagnostics for the EM procedure]{Optimisation
diagnostics at $C=4$. \emph{(a)} Marginal log-likelihood increases
monotonically from four initialisations of $\alpha$. \emph{(b)}
$\widehat\alpha$ converges to $[0.8517,0.8520]$ from starts spanning
$[-0.4,0.6]$, against the true value $0.85$. \emph{(c)} RMSE of
$\widehat\alpha$ and the innovation variance versus $\nseq$, over $16$ seeds,
with an $\nseq^{-1/2}$ reference.}
\label{fig:em-diagnostics}
\end{figure}

\subsection{What does EM recover?}

The parameters are not recovered equally well. Table~\ref{tab:em-recovery}
shows a clear hierarchy: correlation is easiest, innovation variance is next,
and higher-order shape is substantially harder.

\begin{center}
\begin{minipage}{\linewidth}\centering
\captionof{table}[What EM recovers from noised sequences]{Estimates after
$61$ EM iterations on noised Laplace chains
($\alpha=0.85$, $\ivar=0.2775$, excess kurtosis $3.0$, $\nsites=32$).
No run recorded a monotonicity violation.}
\label{tab:em-recovery}
\small
\begin{tabular}{@{}ccccc ccc@{}}
\toprule
$C$ & $\nseq$ & $\widehat\alpha$ & $\widehat\ivar$ & $\widehat\kappa_4$
& $|\Delta\alpha|/\alpha$ & $|\Delta\ivar|/\ivar$ & $|\Delta\kappa_4|/\kappa_4$ \\
\midrule
$4$ & $128$  & $0.8649$ & $0.2686$ & $1.61$ & $1.8\%$  & $3.2\%$ & $46\%$ \\
$4$ & $512$  & $0.8518$ & $0.2844$ & $2.27$ & $0.2\%$  & $2.5\%$ & $24\%$ \\
$4$ & $2048$ & $0.8465$ & $0.2818$ & $1.82$ & $0.4\%$  & $1.5\%$ & $39\%$ \\
\addlinespace
$8$ & $128$  & $0.8294$ & $0.2824$ & $3.18$ & $2.4\%$  & $1.8\%$ & $6\%$ \\
$8$ & $512$  & $0.8503$ & $0.2694$ & $2.26$ & $0.04\%$ & $2.9\%$ & $25\%$ \\
$8$ & $2048$ & $0.8507$ & $0.2793$ & $2.72$ & $0.08\%$ & $0.7\%$ & $9\%$ \\
\bottomrule
\end{tabular}
\end{minipage}
\end{center}

At $C=8$ and $\nseq\geq512$, $\alpha$ is recovered to below $0.1\%$ relative
error, while the innovation variance remains at the percent level and the
fourth cumulant at the tens-of-percent level. Each row is a single fit, and
the shape column is visibly not monotone in $\nseq$ at these budgets --- the
seed-to-seed spread of the fourth cumulant is of the same order as its error,
so these rows demonstrate the recovery hierarchy, not a convergence law for
the shape. The ordering is consistent with
the OU channel: the fourth cumulant is attenuated as $e^{-4t}$, whereas the
second moment is preserved by the variance-matched construction. Higher-order
shape is therefore much less visible in the observations. Throughout,
``recovery'' means what these experiments show: numerical agreement over the
studied regime, with no identifiability or consistency theorem claimed for
the mixture family.

The same conclusion appears directly in the fitted density.

\begin{figure}[t]\centering
\includegraphics[width=\linewidth]{fig_em_innovation.pdf}
\caption[What EM recovers from noised sequences]{Recovery of the innovation
law. \emph{Left:} fitted $C=8$ Gaussian-mixture innovation densities for three
data budgets against the Laplace truth. The main discrepancies occur at the
cusp and in the tails. \emph{Right:} relative errors in $\alpha$, innovation
variance and fourth cumulant versus $\nseq$, showing the hierarchy from
correlation to scale to higher-order shape.}
\label{fig:em-innovation}
\end{figure}

A finite Gaussian mixture necessarily rounds the Laplace cusp, so the remaining
shape error reflects both information lost through the OU channel and
approximation by the mixture family. For the developments that follow, the
estimated kernel is therefore most reliable in its correlation structure,
reasonably accurate in scale, and substantially less informative about
higher-order tail shape.

